\chapter{Discusión}

\emph{Capítulo pendiente de desarrollo.} La discusión completa de este trabajo depende de los resultados de validación descritos como pendientes en la Sección~3.5 y en la Sección~4.4: la reevaluación heurística del prototipo y la prueba de usabilidad con usuarios reales sobre la aplicación completa. Sin esos datos, no corresponde todavía discutir si el prototipo mejora efectivamente la comprensión, la confianza o la percepción de justicia de las familias, que es la pregunta de investigación central de esta memoria (Capítulo~1).

Lo que sí puede discutirse en este punto, por estar basado en resultados ya obtenidos y verificables (Capítulo~4), es lo siguiente.

\section{Por qué el sitio oficial cerró solo dos de las seis brechas prioritarias}

La comparación entre el prototipo y el estado actual del sitio oficial del SAE (Sección~4.3) muestra un patrón que vale la pena anotar como hipótesis a discutir con más evidencia una vez completada la validación: las brechas que el sitio oficial cerró de forma independiente —búsqueda y encontrabilidad, e inclusión parcial mediante el indicador de Programa de Integración Escolar— corresponden a mejoras principalmente informativas y de bajo acoplamiento interactivo, mientras que las que permanecen abiertas —transparencia algorítmica contextualizada, controles interactivos e interoperabilidad con ClaveÚnica— exigen precisamente el tipo de intervención que la literatura revisada en el Capítulo~2 identifica como más efectiva pero también más costosa de implementar: explicaciones locales personalizadas por caso, controles interactivos con retroalimentación inmediata \citep{kim2021recommender,feddersen2024forecasting}, e integración de sistemas de autenticación externos. Esta observación es preliminar y debe tratarse con cautela: compara un diagnóstico de marzo de 2026 con una revisión no sistemática del sitio realizada en julio de 2026, y no controla por otros cambios que el Ministerio de Educación pueda haber implementado fuera de las páginas revisadas.

\section{Qué preguntas deberá responder la discusión completa}

Una vez ejecutado el plan de validación de la Sección~3.5, este capítulo deberá abordar, como mínimo:

\begin{enumerate}
    \item Si el puntaje heurístico del prototipo, reevaluado con el mismo instrumento SISIB del diagnóstico original, efectivamente cierra las brechas críticas identificadas (transparencia, inclusión, búsqueda, audiovisualidad, interoperabilidad), y en qué magnitud respecto del \emph{baseline} de 51\,\%/61\,\%.
    \item Si los resultados de la prueba de usabilidad sobre el prototipo completo son consistentes con los del estudio previo mencionado en la Sección~3.2 —realizado sobre un artefacto más acotado—, en particular respecto de si la explicabilidad contextualizada mejora la comprensión y la confianza sin necesariamente resolver la percepción de justicia del resultado, tal como sugiere la literatura sobre gestión de expectativas \citep{springerwhittaker2019,springerwhittaker2020}.
    \item Qué brechas, si las hubiera, persisten pese al rediseño, y si esas brechas son atribuibles a limitaciones de diseño, a limitaciones inherentes a un prototipo de datos ficticios, o a aspectos del problema que exceden el alcance de una intervención de interfaz (por ejemplo, la interoperabilidad real con ClaveÚnica, que requiere integración de backend no simulable en un prototipo estático).
    \item Qué relación existe entre los hallazgos obtenidos y las condiciones de un despliegue real en producción, dado que la literatura revisada en el Capítulo~2 advierte explícitamente sobre la escasez de estudios de campo frente a estudios de laboratorio.
\end{enumerate}

\section{Limitaciones ya identificables}

Aun sin los resultados de validación, es posible anticipar limitaciones metodológicas que la discusión completa deberá considerar:

\begin{itemize}
    \item \textbf{Datos ficticios}: el catálogo de colegios y los perfiles de usuario del prototipo son ficticios pero representativos del contexto chileno; no se han usado datos reales de postulantes ni de establecimientos.
    \item \textbf{Persona única}: el diseño se validó conceptualmente contra un solo arquetipo de usuario (Daniela González); otros perfiles de apoderados —por ejemplo, familias de estudiantes con necesidades educativas especiales o usuarios migrantes— podrían revelar necesidades de diseño no cubiertas por este trabajo.
    \item \textbf{Alcance del algoritmo}: el prototipo no modifica ni audita el algoritmo de asignación real del SAE, cuyo código fuente no es público; el trabajo se limita a la capa de comunicación de sus resultados.
    \item \textbf{Comparación con el sitio oficial}: la revisión de la Sección~4.3 es una inspección puntual, no un nuevo diagnóstico heurístico completo con el instrumento SISIB; su alcance es más limitado que el del \emph{baseline} original.
\end{itemize}
